% Reconstructed text from V5 PDF page 6; non-native source
\par\smallskip\noindent\rule{\textwidth}{0.4pt}\par{\sffamily\footnotesize Original V5 page 6 -- recovered text}\par\begin{flushleft}\footnotesize\ttfamily
\par\smallskip
\par\smallskip
\noindent Diagnostics and Serving Infrastructure for MAS 6\par
\par\smallskip
\par\smallskip
\par\smallskip
\noindent Sparse activation. Only a fraction of agents invoke the LLM at any moment: a 1,000-agent\par
\noindent simulation may have fewer than 10 actively generating tokens. Under such workloads, keeping\par
\noindent all agent state on the fastest memory tier can waste scarce accelerator memory. ScaleSim’s\par
\noindent invocation-distance metric quantifies the temporal gap before an agent’s next call, enabling\par
\noindent memory-tier demotion for inactive agents.\par
\par\smallskip
\noindent Invocation distance. Many workloads have structured activation patterns such as sequential\par
\noindent queues, event-driven messaging, and dependency graphs. The temporal gap before an agent’s\par
\noindent next call, termed invocation distance, can be estimated from application semantics in structured\par
\noindent workloads but is not generally visible to request-level schedulers. This motivates workflow-\par
\noindent aware caching approaches such as KVFlow, which predicts steps-to-execution for each cached\par
\noindent prefix.\par
\par\smallskip
\noindent Inter-agent memory sharing. Agents may share context such as system prompts, environ-\par
\noindent ment descriptions, and tool schemas. General serving systems already reuse such prefixes:\par
\noindent SGLang’s RadixAttention and vLLM’s automatic prefix caching both exploit shared token se-\par
\noindent quences. The additional opportunity in agentic systems is to use application-level structure\par
\noindent (dependency graphs, message queues, or execution order) to guide cache retention decisions\par
\noindent beyond request-level recency or frequency heuristics.\par
\noindent  The relative importance of these properties varies by application domain. Simulation work-\par
\noindent loads (ScaleSim’s target) exhibit strong sparse activation with predictable invocation patterns.\par
\noindent Coding agent workflows (Continuum’s target) have long tool-call pauses with uncertain du-\par
\noindent rations. Multi-model deployments (DroidSpeak’s target) add cross-architecture cache transfer\par
\noindent challenges. The infrastructure systems reviewed in Section 5 address different subsets of these\par
\noindent workload properties, and the optimal KV-cache management strategy depends on which prop-\par
\noindent erties dominate in a given deployment.\par
\par\smallskip
\noindent 2.3 Runtime Layers\par
\noindent We focus on two runtime layers that are less mature than orchestration:\par
\par\smallskip
\noindent Diagnostics layer. This layer supports failure localization (identifying which agent or step\par
\noindent caused a failure), observability and tracing (continuously capturing execution evidence), and\par
\noindent safety intervention (actively monitoring and preventing harmful actions). The diagnostic layer\par
\noindent consumes execution state — traces, tool outputs, intermediate reasoning — to produce reliabil-\par
\noindent ity outcomes: attribution reports, anomaly alerts, and safety blocks. The surveyed methods vary\par
\noindent dramatically in their resource requirements: FAMAS requires 20 environment replays averag-\par
\noindent ing 105 minutes per task, while MASPrism requires only two deterministic prefill passes at 2.66\par
\noindent seconds on a 0.6B model. This three-orders-of-magnitude cost range has direct implications for\par
\noindent cross-layer design: methods at the high end of the cost spectrum may compete with serving\par
\noindent resources, while methods at the low end operate well within available headroom.\par
\par\smallskip
\noindent Infrastructure layer. This layer manages serving-time resources: KV-cache memory, com-\par
\noindent pute scheduling, adapter placement, and power management for multi-agent workloads. The\par
\noindent infrastructure layer consumes serving state — KV blocks, adapters, request metadata — and\par
\noindent optimizes for throughput, latency, and resource utilization. Among the 22 Tier-1 infrastructure\par
\end{flushleft}
